\documentclass[11pt,a4paper]{article}

\usepackage[utf8]{inputenc}
\usepackage[T1]{fontenc}
\usepackage{lmodern}
\usepackage{amsmath,amssymb,amsthm}
\usepackage{mathtools}
\usepackage{hyperref}
\usepackage[margin=1in]{geometry}

% Theorem environments
\newtheorem{theorem}{Theorem}[section]
\newtheorem{proposition}[theorem]{Proposition}
\newtheorem{corollary}[theorem]{Corollary}
\newtheorem{definition}[theorem]{Definition}
\newtheorem{example}[theorem]{Example}
\newtheorem{remark}[theorem]{Remark}
\newtheorem{problem}[theorem]{Problem}

\title{\textbf{Stratified Gauge Theory:\\Invariance and Conservation on Singular Spaces}\\[0.5em]
\large Extending Noether's Theorem to Stratified Principal Bundles}

\author{Andrew H. Bond\\
\textit{Department of Computer Engineering}\\
\textit{San Jos\'e State University}\\
\texttt{andrew.bond@sjsu.edu}}

\date{December 2025}

\begin{document}

\maketitle

\begin{abstract}
We develop a gauge-theoretic framework compatible with stratified spaces, resolving the
tension between (1) the smooth manifold structure required by standard Noether conservation and (2) the stratified structure required for moral discontinuities, hard constraints,
and regime changes. The key insight is that gauge structure and Noether conservation hold
\textit{stratum-wise}, while boundary behavior is governed by \textit{matching conditions} that constrain
how trajectories cross between strata. We define stratified principal bundles, stratified
connections, and prove a \textbf{Stratified Noether Theorem}: the alignment current is conserved within each stratum and satisfies explicit jump conditions at stratum boundaries.
For the special case of constraint boundaries (where evaluation jumps to $-\infty$), we prove
that admissible trajectories cannot cross the boundary, recovering the No Escape result as a
consequence of the variational principle rather than an additional assumption. The framework unifies Stratified Geometric Ethics with gauge-theoretic invariance while maintaining
mathematical rigor.
\end{abstract}

\medskip
\noindent\textbf{Keywords:} stratified spaces, gauge theory, Noether's theorem, principal bundles, conservation laws, singular spaces, jump conditions, variational principles, geometric ethics

\tableofcontents

%==============================================================================
\section{Introduction: The Smoothness Problem}
%==============================================================================

\subsection{The Tension}

Two papers in this series make apparently conflicting claims (see~\cite{bond2026ge} for the unified development of the Geometric Ethics framework within which this work is situated):

\begin{enumerate}
    \item \textbf{Stratified Geometric Ethics (SGE):} The space of morally relevant configurations must be modeled as a \textit{stratified space}---a union of smooth manifolds of varying dimensions---to represent:
    \begin{itemize}
        \item Hard vetoes (regions where $\Sigma = -\infty$)
        \item Lexical priorities (infinite cost ratios)
        \item Genuine dilemmas (singular configurations)
        \item Regime changes (discontinuous shifts in applicable rules)
    \end{itemize}

    \item \textbf{Gauge Theory of Invariant Control:} Alignment is a conserved Noether current, which requires:
    \begin{itemize}
        \item Smooth configuration manifold
        \item Lie group symmetry
        \item Smooth Lagrangian
        \item Euler-Lagrange dynamics
    \end{itemize}
\end{enumerate}

\begin{problem}[The Smoothness Problem]
Noether's theorem requires smoothness. Stratified spaces have singular boundaries where smoothness fails. How can conservation hold on stratified spaces?
\end{problem}

\subsection{The Resolution Strategy}

We resolve the tension by observing:

\begin{enumerate}
    \item Each \textit{stratum} is a smooth manifold. Gauge theory and Noether's theorem apply within each stratum.

    \item Stratum \textit{boundaries} are where smoothness fails. We need additional structure to handle boundary crossing.

    \item The \textit{matching conditions} at boundaries determine whether trajectories can cross and how conserved quantities behave at crossings.

    \item For \textit{constraint boundaries} (where $\Sigma$ jumps to $-\infty$), the variational principle itself forbids crossing---no additional ``enforcement'' is needed.
\end{enumerate}

\subsection{Contributions}

\begin{enumerate}
    \item \textbf{Definition:} Stratified principal bundles and stratified connections (Section~\ref{sec:bundles})

    \item \textbf{Theorem:} Stratum-wise Noether conservation (Section~\ref{sec:noether})

    \item \textbf{Theorem:} Jump conditions at boundaries (Section~\ref{sec:boundary})

    \item \textbf{Theorem:} Constraint boundaries are impassable (Section~\ref{sec:full-noether})

    \item \textbf{Application:} Regime-dependent control with conservation (Section~\ref{sec:reconciliation})
\end{enumerate}

%==============================================================================
\section{Background: Stratified Spaces}\label{sec:background}
%==============================================================================

We recall the essential definitions. See Pflaum~\cite{pflaum} or Mather~\cite{mather} for full treatments.

\subsection{Whitney Stratified Spaces}

\begin{definition}[Stratified Space]
A \textbf{stratified space} is a triple $(\mathcal{M}, \{M_i\}_{i \in I}, \preceq)$ where:
\begin{itemize}
    \item $\mathcal{M}$ is a paracompact Hausdorff topological space
    \item $\{M_i\}_{i \in I}$ is a locally finite partition of $\mathcal{M}$ into connected smooth manifolds (called \textbf{strata})
    \item $\preceq$ is a partial order on $I$ satisfying the \textbf{frontier condition}:
\end{itemize}
\[
    M_i \cap \mathrm{cl}(M_j) \neq \emptyset \implies i \preceq j
\]
\end{definition}

\begin{definition}[Whitney Condition B]
The stratification satisfies \textbf{Whitney's condition (B)} if: for sequences $\{y_n\} \subset M_j$, $\{x_n\} \subset M_i$ with $y_n, x_n \to x \in M_i$ $(i \prec j)$, if the secant lines $\ell_n = \overline{x_n y_n}$ converge to $\ell$ and the tangent spaces $T_{y_n} M_j$ converge to $\tau$, then $\ell \subseteq \tau$.
\end{definition}

\begin{remark}
Whitney (B) ensures that strata fit together ``nicely'': tangent spaces of higher strata contain the limiting directions from lower strata. This is essential for well-defined limits at boundaries.
\end{remark}

\subsection{Stratified Maps and Dynamics}

\begin{definition}[Stratified Map]
A continuous map $f : \mathcal{M} \to \mathcal{M}'$ between stratified spaces is \textbf{stratified} if for each stratum $M_i$ of $\mathcal{M}$, $f(M_i)$ is contained in a single stratum of $\mathcal{M}'$ and $f|_{M_i}$ is smooth.
\end{definition}

\begin{definition}[Stratified Vector Field]
A \textbf{stratified vector field} on $\mathcal{M}$ is a continuous section $X : \mathcal{M} \to T\mathcal{M}$ such that $X|_{M_i}$ is a smooth vector field on $M_i$ for each stratum.
\end{definition}

\begin{definition}[Stratified Curve]
A continuous curve $\gamma : [0, T] \to \mathcal{M}$ is \textbf{stratified} if:
\begin{enumerate}
    \item[(i)] There exists a finite partition $0 = t_0 < t_1 < \cdots < t_k = T$
    \item[(ii)] On each interval $(t_{j-1}, t_j)$, the curve lies in a single stratum and is smooth
    \item[(iii)] At each $t_j$, the curve may transition between strata
\end{enumerate}
\end{definition}

%==============================================================================
\section{Stratified Principal Bundles}\label{sec:bundles}
%==============================================================================

\subsection{Definition}

\begin{definition}[Stratified Principal Bundle]
A \textbf{stratified principal $\mathcal{G}$-bundle} over a stratified base $(\mathcal{M}, \{M_i\}, \preceq)$ is a tuple $(\mathcal{P}, \mathcal{M}, \pi, \mathcal{G}, \{P_i\})$ where:
\begin{enumerate}
    \item[\textup{(SPB1)}] $\mathcal{P}$ is a stratified space with strata $\{P_i\}_{i \in I}$
    \item[\textup{(SPB2)}] $\pi : \mathcal{P} \to \mathcal{M}$ is a continuous surjection
    \item[\textup{(SPB3)}] $\mathcal{G}$ is a Lie group acting continuously on $\mathcal{P}$ from the right
    \item[\textup{(SPB4)}] For each stratum $M_i$, the restriction $\pi_i : P_i \to M_i$ is a (smooth) principal $\mathcal{G}$-bundle
    \item[\textup{(SPB5)}] $\pi$ respects stratification: $\pi(P_i) = M_i$
    \item[\textup{(SPB6)}] The $\mathcal{G}$-action respects stratification: $P_i \cdot g = P_i$ for all $g \in \mathcal{G}$
\end{enumerate}
\end{definition}

\begin{proposition}[Stratum-wise Bundle Structure]
Each restriction $\pi_i : P_i \to M_i$ has the standard structure of a smooth principal bundle: free $\mathcal{G}$-action, local trivializations, transition functions, etc.
\end{proposition}

\begin{proof}
Immediate from (SPB4).
\end{proof}

\subsection{Stratified Connections}

\begin{definition}[Stratified Connection]
A \textbf{stratified connection} on a stratified principal bundle $(\mathcal{P}, \mathcal{M}, \pi, \mathcal{G})$ is a collection $\boldsymbol{\omega} = \{\omega_i\}_{i \in I}$ where:
\begin{enumerate}
    \item[(i)] Each $\omega_i$ is a connection 1-form on the smooth bundle $P_i \to M_i$
    \item[(ii)] The connections satisfy \textbf{compatibility conditions} at boundaries (see below)
\end{enumerate}
\end{definition}

\begin{definition}[Boundary Compatibility]
Let $M_i \prec M_j$ (so $M_i \subset \mathrm{cl}(M_j)$). The connections $\omega_i$ and $\omega_j$ are \textbf{compatible} if for any sequence $p_n \in P_j$ with $\pi(p_n) \to m \in M_i$ and any sequence of tangent vectors $v_n \in T_{p_n} P_j$ converging to $v \in T_p P_i$ (where $p \in \pi^{-1}(m)$), we have:
\[
    \lim_{n \to \infty} \omega_j(v_n) = \omega_i(v)
\]
whenever the limit exists.
\end{definition}

\begin{remark}
Compatibility ensures that the connection forms ``match up'' at stratum boundaries. This is non-trivial: the connections live on different bundles, and the limit involves both spatial convergence ($\pi(p_n) \to m$) and tangent space convergence.
\end{remark}

\subsection{Stratified Curvature}

\begin{definition}[Stratified Curvature]
The \textbf{stratified curvature} is the collection $\boldsymbol{\Omega} = \{\Omega_i\}_{i \in I}$ where each $\Omega_i$ is the curvature 2-form of $\omega_i$:
\[
    \Omega_i = d\omega_i + \frac{1}{2}[\omega_i, \omega_i]
\]
\end{definition}

\begin{definition}[Stratum-wise Flatness]
A stratified connection is \textbf{stratum-wise flat} if $\Omega_i = 0$ for all $i \in I$.
\end{definition}

\begin{theorem}[Canonicalization Gives Stratum-wise Flat Connection]
Let $(\mathcal{P}, \mathcal{M}, \pi, \mathcal{G})$ be a stratified principal bundle. Suppose there exists a \textbf{stratified canonicalizer} $\kappa : \mathcal{P} \to \mathcal{P}$ such that:
\begin{enumerate}
    \item[(i)] $\kappa$ is stratified (maps each $P_i$ to itself)
    \item[(ii)] $\kappa$ restricted to each $P_i$ satisfies the canonicalizer axioms
\end{enumerate}
Then $\kappa$ induces a stratum-wise flat stratified connection.
\end{theorem}

\begin{proof}
On each stratum $P_i$, the canonicalizer induces a global section $\sigma_i : M_i \to P_i$. This section determines a flat connection $\omega_i$ with $\Omega_i = 0$ by the standard argument (trivialization via the section). The collection $\{\omega_i\}$ forms a stratum-wise flat stratified connection.
\end{proof}

%==============================================================================
\section{Stratified Noether Theorem}\label{sec:noether}
%==============================================================================

\subsection{Stratum-wise Lagrangian Dynamics}

\begin{definition}[Stratified Lagrangian]
A \textbf{stratified Lagrangian} on a stratified principal bundle is a collection $\mathbf{L} = \{L_i\}_{i \in I}$ where each $L_i : TP_i \to \mathbb{R}$ is a smooth Lagrangian on the $i$-th stratum.
\end{definition}

\begin{definition}[Gauge-Invariant Stratified Lagrangian]
A stratified Lagrangian is \textbf{gauge-invariant} if each $L_i$ is $\mathcal{G}$-invariant:
\[
    L_i(p \cdot g, (R_g)_* v) = L_i(p, v) \quad \forall g \in \mathcal{G},\; p \in P_i,\; v \in T_p P_i
\]
\end{definition}

\begin{example}[Decision Lagrangian]
On each stratum:
\[
    L_i(p, \dot{p}) = \frac{1}{2} g_{M_i}(\check{\dot{p}}, \check{\dot{p}}) - U_i(\pi(p)) - \frac{\lambda}{2} \|\omega_i(\dot{p})\|^2
\]
where $g_{M_i}$ is a Riemannian metric on $M_i$, $U_i : M_i \to \mathbb{R}$ is a potential, and $\omega_i$ is the connection form.
\end{example}

\subsection{Stratum-Interior Dynamics}

\begin{definition}[Interior Trajectory]
A curve $\gamma : (a, b) \to P_i$ lying entirely within a single stratum is an \textbf{interior trajectory}.
\end{definition}

\begin{theorem}[Stratum-Interior Euler-Lagrange Equations]
An interior trajectory $\gamma : (a, b) \to P_i$ is a critical point of the action $S_i[\gamma] = \int_a^b L_i(\gamma, \dot{\gamma})\,dt$ if and only if it satisfies the Euler-Lagrange equations:
\[
    \frac{d}{dt}\frac{\partial L_i}{\partial \dot{p}} - \frac{\partial L_i}{\partial p} = 0
\]
\end{theorem}

\begin{proof}
Standard calculus of variations on the smooth manifold $P_i$.
\end{proof}

\subsection{The Stratum-Interior Noether Theorem}

\begin{theorem}[Stratum-Interior Noether Conservation]
Let $\mathbf{L}$ be a gauge-invariant stratified Lagrangian. For any interior trajectory $\gamma : (a, b) \to P_i$ satisfying the Euler-Lagrange equations, the \textbf{alignment current}
\[
    J_i := \lambda \omega_i(\dot{\gamma}) \in \mathfrak{g}
\]
is conserved: $\frac{d}{dt} J_i = 0$ on $(a, b)$.
\end{theorem}

\begin{proof}
This is the standard Noether theorem applied to the smooth bundle $P_i \to M_i$ with Lie group symmetry $\mathcal{G}$. The Lagrangian $L_i$ is $\mathcal{G}$-invariant by assumption. For each $\xi \in \mathfrak{g}$, the infinitesimal generator $\xi_{P_i}$ gives a Noether charge $Q_\xi = \langle \partial L_i / \partial \dot{p}, \xi_{P_i} \rangle$. Computing as in the Gauge Theory paper, $Q_\xi = \lambda \langle \omega_i(\dot{\gamma}), \xi \rangle$, so $J_i = \lambda \omega_i(\dot{\gamma})$ is the full Noether current. Conservation follows from the Euler-Lagrange equations.
\end{proof}

\begin{corollary}[Stratum-Interior Horizontal Preservation]
If $\gamma$ begins horizontal ($\omega_i(\dot{\gamma}(a)) = 0$), it remains horizontal throughout $(a, b)$.
\end{corollary}

\begin{proof}
$J_i(t) = J_i(a) = 0$ for all $t \in (a, b)$.
\end{proof}

%==============================================================================
\section{Boundary Behavior and Jump Conditions}\label{sec:boundary}
%==============================================================================

We now address what happens when a trajectory approaches or crosses a stratum boundary.

\subsection{Types of Boundaries}

\begin{definition}[Boundary Types]
Let $M_i \prec M_j$ be adjacent strata ($M_i \subset \mathrm{cl}(M_j) \setminus M_j$). The boundary is:
\begin{enumerate}
    \item \textbf{Regular} if the potentials $U_i$ and $U_j$ extend continuously to the common boundary
    \item \textbf{Discontinuous} if $\lim_{m \to \partial} U_j(m) \neq U_i|_\partial$ (finite jump)
    \item \textbf{Singular (constraint boundary)} if $U_i = +\infty$ on $M_i$ (equivalently, $\Sigma_i = -\infty$)
\end{enumerate}
\end{definition}

\subsection{Regular Boundaries: Matching Conditions}

\begin{theorem}[Regular Boundary Crossing]
At a regular boundary between strata $M_j$ and $M_i$ $(i \prec j)$, a trajectory $\gamma$ crossing from $M_j$ to $M_i$ at time $t^*$ satisfies:
\begin{enumerate}
    \item[(i)] \textbf{Position continuity:} $\lim_{t \to t^{*-}} \gamma(t) = \lim_{t \to t^{*+}} \gamma(t)$
    \item[(ii)] \textbf{Momentum matching:} The generalized momenta match across the boundary (Weierstrass-Erdmann conditions)
    \item[(iii)] \textbf{Current matching:} If the connections are compatible, $\lim_{t \to t^{*-}} J_j(t) = \lim_{t \to t^{*+}} J_i(t)$
\end{enumerate}
\end{theorem}

\begin{proof}
(i) Continuity is required for the trajectory to be a valid curve in $\mathcal{M}$.

(ii) The Weierstrass-Erdmann corner conditions for variational problems with discontinuities require:
\[
    \lim_{t \to t^{*-}} \frac{\partial L_j}{\partial \dot{p}} = \lim_{t \to t^{*+}} \frac{\partial L_i}{\partial \dot{p}}
\]
when the Lagrangians match continuously at the boundary.

(iii) The momentum $\partial L / \partial \dot{p}$ includes the term $\lambda \omega(\dot{p})$ from the gauge-fixing term. If connections are compatible (boundary compatibility condition), the connection forms match in the limit, so the Noether currents match.
\end{proof}

\begin{corollary}[Conservation Across Regular Boundaries]
At a regular boundary with compatible connections, the alignment current is continuous: there is no jump in $J$.
\end{corollary}

\subsection{Discontinuous Boundaries: Jump Conditions}

\begin{theorem}[Discontinuous Boundary Jump]
At a discontinuous boundary where $U_j \to U_j^-$ and $U_i = U_i^+$ with $\Delta U := U_i^+ - U_j^- \neq 0$, the trajectory experiences:
\begin{enumerate}
    \item[(i)] \textbf{Velocity jump:} The kinetic energy changes by $\Delta K = -\Delta U$ (energy conservation)
    \item[(ii)] \textbf{Current jump:} The alignment current may jump: $\Delta J = J_i^+ - J_j^- \neq 0$ in general
\end{enumerate}

The jump in $J$ is determined by how the velocity jump projects onto the fiber direction.
\end{theorem}

\begin{proof}
Energy conservation across the discontinuity gives $K_j^- + U_j^- = K_i^+ + U_i^+$, hence $\Delta K = -\Delta U$.

The Noether current $J = \lambda \omega(\dot{p})$ depends on $\dot{p}$. If the velocity jumps, and the jump has a component in the fiber direction, then $J$ jumps. Specifically:
\[
    \Delta J = \lambda \omega_i(\dot{p}_i^+) - \lambda \omega_j(\dot{p}_j^-)
\]
Even with compatible connections, $\omega_i(\dot{p}_i^+) \neq \omega_j(\dot{p}_j^-)$ if the velocity changes.
\end{proof}

\begin{remark}
At discontinuous boundaries, the Noether current is \textit{not} conserved across the boundary. Conservation holds \textit{within} each stratum, but there can be ``injection'' or ``absorption'' of alignment current at regime changes.
\end{remark}

\subsection{Singular (Constraint) Boundaries: Impassability}

This is the crucial case for the No Escape theorem.

\begin{theorem}[Constraint Boundaries are Impassable]\label{thm:impassable}
Let $M_i$ be a constraint stratum with $U_i = +\infty$ (equivalently, $\Sigma_i = -\infty$). Then no finite-action trajectory can enter $M_i$.
\end{theorem}

\begin{proof}
Suppose $\gamma : [0, T] \to \mathcal{P}$ is a trajectory that enters $M_i$ at some time $t^* \in (0, T)$. The action is:
\[
    S[\gamma] = \int_0^T L(\gamma, \dot{\gamma})\,dt
\]
For $t > t^*$, the trajectory is in $P_i$, where the Lagrangian includes the term $-U_i(\pi(\gamma(t))) = -\infty$.
Therefore:
\[
    S[\gamma] = \int_0^{t^*} L_j(\gamma, \dot{\gamma})\,dt + \int_{t^*}^T L_i(\gamma, \dot{\gamma})\,dt = \text{finite} + (-\infty) = -\infty
\]
But critical points of the action are trajectories with finite action (the Euler-Lagrange equations are derived by setting the variation to zero, which requires finite values). A trajectory with $S = -\infty$ is not a critical point in any meaningful sense.

More directly: trajectories are chosen to minimize/extremize the action. A trajectory entering the constraint region has $S = -\infty$, which is lower than any finite-action trajectory. But the ``infimum'' $-\infty$ is not achieved by any actual trajectory---it represents the limit of trajectories that approach but never enter $M_i$. The optimal trajectories are those that stay outside $M_i$.
\end{proof}

\begin{corollary}[Variational No Escape]
If the constraint region $\mathcal{C}^c = \bigcup_{i: U_i = +\infty} M_i$ is defined by $U = +\infty$, then all Euler-Lagrange trajectories remain in the admissible region $\mathcal{C} = \mathcal{M} \setminus \mathcal{C}^c$.
\end{corollary}

\begin{remark}[No Escape as Variational Consequence]
This is a key result. In the smooth gauge theory paper, the No Escape theorem was stated as a consequence of structural containment. Here we see it more directly: \textit{constraint violation is forbidden by the variational principle itself}. The Lagrangian encodes constraints as infinite potential barriers; the dynamics cannot cross them.
\end{remark}

%==============================================================================
\section{The Stratified Noether Theorem: Full Statement}\label{sec:full-noether}
%==============================================================================

We now combine the results into a comprehensive statement.

\begin{theorem}[Stratified Noether Theorem]\label{thm:stratified-noether}
Let $(\mathcal{P}, \mathcal{M}, \pi, \mathcal{G})$ be a stratified principal bundle with gauge-invariant stratified Lagrangian $\mathbf{L} = \{L_i\}$ and stratified connection $\boldsymbol{\omega} = \{\omega_i\}$. Let $\gamma : [0, T] \to \mathcal{P}$ be a piecewise-smooth trajectory that:
\begin{itemize}
    \item Lies in the admissible region (avoids constraint strata)
    \item Satisfies Euler-Lagrange equations within each stratum
    \item Satisfies Weierstrass-Erdmann conditions at boundary crossings
\end{itemize}
Define the alignment current $J(t) := \lambda \omega_{i(t)}(\dot{\gamma}(t))$ where $i(t)$ is the stratum containing $\gamma(t)$. Then:
\begin{enumerate}
    \item[(a)] \textbf{Stratum-interior conservation:} $\frac{d}{dt} J = 0$ on intervals where $\gamma$ lies in a single stratum
    \item[(b)] \textbf{Regular boundary:} At crossings between strata with continuous potentials and compatible connections, $J$ is continuous
    \item[(c)] \textbf{Discontinuous boundary:} At crossings with potential jumps $\Delta U$, the current may jump by an amount determined by the velocity change
    \item[(d)] \textbf{Constraint boundary:} Trajectories cannot reach strata with $U = +\infty$
    \item[(e)] \textbf{Horizontal preservation (within strata):} If $J(t_0) = 0$ for some $t_0$ in the interior of a stratum, then $J(t) = 0$ for all $t$ in that stratum-interior segment
    \item[(f)] \textbf{Global horizontal preservation (regular case):} If all boundaries are regular with compatible connections and $J(0) = 0$, then $J(t) = 0$ for all $t \in [0, T]$
\end{enumerate}
\end{theorem}

\begin{proof}
(a) Theorem~4.6. (b) Theorem~5.2. (c) Theorem~5.4. (d) Theorem~\ref{thm:impassable}. (e) Corollary of (a): within a stratum, Noether conservation applies. (f) Combine (a), (b), and (e): $J$ is conserved within strata and continuous across regular boundaries, so globally $J(t) = J(0) = 0$.
\end{proof}

%==============================================================================
\section{Reconciliation with SGE}\label{sec:reconciliation}
%==============================================================================

\subsection{The Original Tension, Resolved}

The apparent conflict was:
\begin{itemize}
    \item SGE requires stratification for discontinuities
    \item Noether requires smoothness for conservation
\end{itemize}

The resolution:
\begin{itemize}
    \item Noether conservation holds \textit{within} each stratum (where everything is smooth)
    \item Boundary behavior is governed by matching conditions (regular case) or impassability (constraint case)
    \item The stratified structure is essential for modeling hard constraints, which appear as infinite potential barriers
\end{itemize}

\subsection{Hard Vetoes as Infinite Barriers}

\begin{proposition}[SGE Hard Veto = Infinite Potential]
In SGE, a hard veto is a region where $\Sigma = -\infty$ (absolutely forbidden). In the Lagrangian formulation, this corresponds to $U = -\Sigma = +\infty$. By Theorem~\ref{thm:impassable}, Euler-Lagrange trajectories cannot enter hard veto regions.
\end{proposition}

\textbf{Interpretation:} Hard vetoes are not enforced by an external mechanism. They are \textit{built into the dynamics} as infinite potential barriers. The variational principle itself forbids violation.

\subsection{Lexical Priorities as Metric Singularities}

\begin{proposition}[Lexical Priority as Singular Metric]
In SGE, lexical priority of value $v_1$ over $v_2$ means: any decrease in $v_1$ dominates any increase in $v_2$. This is modeled by a family of metrics $g_\varepsilon$ with $g_\varepsilon^{11} \sim 1/\varepsilon^2$ as $\varepsilon \to 0$.

In the stratified framework, the limit $\varepsilon \to 0$ is a stratum boundary where the metric becomes singular.
\end{proposition}

\textbf{Interpretation:} Lexical priorities correspond to metric singularities at stratum boundaries. Motion in the lexically prior direction becomes infinitely costly, so trajectories avoid it.

\subsection{Genuine Dilemmas as Singular Strata}

\begin{proposition}[Dilemma = Singular Stratum]
In SGE, a genuine dilemma is a configuration from which all exits incur positive moral cost. This corresponds to a low-dimensional stratum $M_i$ (possibly 0-dimensional) from which all paths to higher-dimensional strata pass through regions of high $U$.
\end{proposition}

\textbf{Interpretation:} Dilemmas are topological features of the stratified space. They are not ``solved'' by the dynamics---they represent configurations where the agent is trapped by the geometry.

%==============================================================================
\section{Application: Regime-Dependent Control}\label{sec:application}
%==============================================================================

\subsection{Plasma Control with Regime Transitions}

Tokamak plasmas exhibit distinct operating regimes:
\begin{itemize}
    \item \textbf{L-mode:} Low confinement, standard transport
    \item \textbf{H-mode:} High confinement, edge pedestal
    \item \textbf{ELMy H-mode:} H-mode with edge-localized modes
\end{itemize}

Each regime has different physics and different control objectives. Transitions between regimes are discontinuous.

\subsection{Stratified Model}

\begin{itemize}
    \item \textbf{Strata:} $M_L$ (L-mode), $M_H$ (H-mode), $M_E$ (ELMy H-mode)
    \item \textbf{Grounding:} $\Psi = (f_G, \beta_N, \nu^*, H_{98}, \ldots)$ (dimensionless plasma parameters)
    \item \textbf{Potentials:} $U_L, U_H, U_E$ (different in each regime)
    \item \textbf{Constraints:} $U = +\infty$ for $f_G > 1$ (Greenwald limit), $\beta_N > \beta_{\mathrm{crit}}$ (pressure limit)
\end{itemize}

\subsection{Conservation and Transitions}

\begin{proposition}[Regime-Dependent Control Properties]\leavevmode
\begin{enumerate}
    \item \textbf{Within each regime:} Alignment current $J$ is conserved; horizontal trajectories remain horizontal.
    \item \textbf{At regime transitions:} If the transition is modeled as a discontinuous boundary, $J$ may jump. The controller must re-establish alignment after the transition.
    \item \textbf{At constraint boundaries:} Trajectories cannot violate Greenwald or beta limits; these are variationally forbidden.
\end{enumerate}
\end{proposition}

\subsection{Practical Implications}

\begin{enumerate}
    \item \textbf{Intra-regime control:} Within a regime, BIP-compliant controllers stay BIP-compliant (conservation).
    \item \textbf{Inter-regime control:} At regime transitions, the controller may need recalibration. The jump conditions specify how.
    \item \textbf{Safety:} Hard limits (Greenwald, beta) are never violated, regardless of regime.
\end{enumerate}

%==============================================================================
\section{Conclusion}
%==============================================================================

\subsection{Summary of Results}

\begin{enumerate}
    \item \textbf{Stratified principal bundles} extend gauge theory to stratified spaces.

    \item \textbf{Stratified connections} provide gauge structure on each stratum with matching conditions at boundaries.

    \item \textbf{Stratum-interior Noether conservation} holds: alignment is conserved within each smooth stratum.

    \item \textbf{Boundary behavior} is classified:
    \begin{itemize}
        \item Regular boundaries: current is continuous
        \item Discontinuous boundaries: current may jump
        \item Constraint boundaries: trajectories are blocked
    \end{itemize}

    \item \textbf{No Escape} is a variational consequence: constraint regions with $U = +\infty$ are impassable.
\end{enumerate}

\subsection{Resolution of the Smoothness Problem}

The apparent conflict between stratification (needed for discontinuities) and Noether theory (needing smoothness) is resolved by:
\begin{enumerate}
    \item Applying Noether theory \textit{within} each stratum (where smoothness holds)
    \item Supplementing with matching conditions at boundaries
    \item Recognizing constraint boundaries as infinite barriers in the variational problem
\end{enumerate}

\subsection{The Unified Picture}

The stratified gauge theory framework unifies:
\begin{itemize}
    \item \textbf{SGE's} stratified moral geometry
    \item \textbf{BIP's} gauge invariance requirement
    \item \textbf{Noether conservation} (stratum-wise)
    \item \textbf{No Escape} (as variational impassability)
    \item \textbf{Regime-dependent dynamics} (with explicit jump conditions)
\end{itemize}
into a single mathematically coherent structure. A reference implementation of this framework, including stratified evaluation and tensorial ethics, is available in ErisML/DEME v3.0~\cite{bond2026erisml}.

%==============================================================================
\begin{thebibliography}{9}

\bibitem{pflaum}
M.J.\ Pflaum, \textit{Analytic and Geometric Study of Stratified Spaces}, Springer Lecture Notes in Mathematics 1768, 2001.

\bibitem{mather}
J.\ Mather, ``Notes on topological stability,'' \textit{Bulletin of the American Mathematical Society}, 49(4):475--506, 2012.

\bibitem{arnold}
V.I.\ Arnold, \textit{Mathematical Methods of Classical Mechanics}, 2nd ed., Springer, 1989.

\bibitem{kobayashi}
S.\ Kobayashi and K.\ Nomizu, \textit{Foundations of Differential Geometry}, Wiley, 1963.

\bibitem{bond2026ge}
A.~H.~Bond, \emph{Geometric Ethics: The Mathematical Structure of Moral Reasoning}, v1.14, Department of Computer Engineering, San Jos\'e State University, Spring 2026.

\bibitem{bond2026erisml}
A.~H.~Bond, ``ErisML/DEME v3.0: A Library for Governed, Foundation-Model-Enabled Agents with Tensorial Ethics,'' \url{https://github.com/ahb-sjsu/erisml-lib}, 2026.

\end{thebibliography}

\end{document}
